\documentclass[a4paper,11pt]{article}
\usepackage[utf8]{inputenc}
\usepackage[T1]{fontenc}
\usepackage[french]{babel}
\usepackage[left=2cm,right=2cm,top=2cm,bottom=2cm]{geometry}
\usepackage{xcolor}
\usepackage{hyperref}
\usepackage{fontawesome5}
\usepackage{parskip}

\definecolor{primary}{RGB}{35, 55, 123}

\begin{document}

\begin{center}
    {\Large \textbf{Loïc Aron MBASSI EWOLO}}\\[4pt]
    \small
    \faMapMarker\ Cergy, France \quad
    \faPhone\ (+33) 7 59 52 33 98 \quad
    \faEnvelope\ \href{mailto:loicaron.mbassiewolo@ensea.fr}{loicaron.mbassiewolo@ensea.fr}
\end{center}
\vspace{0.5cm}

\hfill
\begin{minipage}[t]{0.45\textwidth}
    \textbf{KuroKoma Studio}\\
    Service Recrutement / Engineering Team
\end{minipage}

\vspace{1cm}

\textbf{Objet :} Candidature Stage - GenAI Product Engineer (Avril 2026)

Bonjour,

Le défi de la cohérence d'identité dans la génération d'images est le Saint Graal actuel de la GenAI. Construire un outil studio-grade pour le manga ne demande pas seulement de bons modèles, mais un excellent engineering autour. C'est exactement là que je veux intervenir.

Mon profil de Product Engineer correspond à votre besoin d'hybridation entre IA et Dev :

\textbf{1. Python & Streamlit (Build) :}
Vous cherchez quelqu'un à l'aise avec Streamlit ? C'est la technologie que j'ai choisie pour développer le dashboard qui m'a valu la \textbf{2ème place au hackathon IndabaX Africa}. Je sais prototyper rapidement des interfaces data-centric pour rendre l'IA accessible aux utilisateurs finaux.

\textbf{2. Pipelines IA & Workflows (Automate) :}
J'ai conçu un système multi-agents complexe (Google ADK) intégrant du RAG et de la gestion de contexte. Cette expérience de l'orchestration de flux IA me permettra de m'attaquer efficacement à vos pipelines Stable Diffusion / LoRA et d'automatiser les tâches répétitives pour les créateurs.

\textbf{3. Mentalité Shipping (Deliver) :}
Au-delà de l'école (ENSEA/Polytech), j'ai une expérience freelance concrète. J'ai livré des applications en production (React, Laravel), géré du versioning Git et mis en place des environnements Docker. Je comprends qu'un outil interne doit être robuste, documenté et maintenable.

Passionné par l'intersection entre créativité et technologie, je suis disponible dès avril 2026 pour aider KuroKoma à devenir la référence des mangakas.

Cordialement,

\vspace{1cm}

\textbf{Loïc Aron MBASSI EWOLO}

\end{document}